% Generated draft: requires mathematical and source review.
\documentclass[11pt]{article}
\usepackage[margin=25mm]{geometry}
\usepackage{fontspec}
\setmainfont{Latin Modern Roman}
\usepackage{amsmath,amssymb,mathtools,unicode-math}
\setmathfont{Latin Modern Math}
\usepackage{enumitem,needspace,xurl,hyperref}
\setlength{\parindent}{0pt}
\setlength{\parskip}{5pt}
\setlength{\emergencystretch}{3em}
\newcommand{\N}{\mathbb{N}}
\newcommand{\Z}{\mathbb{Z}}
\newcommand{\Q}{\mathbb{Q}}
\newcommand{\R}{\mathbb{R}}
\newcommand{\C}{\mathbb{C}}
\newcommand{\F}{\mathbb{F}}
\newcommand{\A}{\mathbb{A}}
\newcommand{\PP}{\mathbb{P}}
\newcommand{\E}{\mathbb{E}}
\newcommand{\eps}{\varepsilon}
\newcommand{\dd}{\,\mathrm{d}}
\begin{document}
\section*{UnsolvedMath ID: 30004799}
\textbf{Source problem code:} catalogue entry
\subsection{Definitions and mathematical statement}
Let \(v\), \(G_v\), and \(\psi\) be a place, a local group, and an Arthur parameter, respectively, within the scope of the construction cited in the source record. Write
\[
\Pi^{\theta(\psi)}_v(G_v)
\]
for the local packet attached to \(\psi\) by that construction. The supplied record does not specify the permitted groups, the precise definition of an Arthur parameter in this setting, or the packet construction. These are source-dependent data, not unrestricted choices.

To express the auxiliary-choice question, write \(\Pi^{\theta(\psi)}_{v,H}(G_v)\) for this packet when the auxiliary group is \(H\). This notation merely records a choice already present in the source construction; the admissible auxiliary groups are not specified in the supplied record.

The requested assertions are as follows, with all hypotheses and admissibility conditions of the cited construction retained:
\begin{enumerate}
\item For every local place \(v\) and every Arthur parameter \(\psi\) covered by the construction, establish the endoscopic character identities for \(\Pi^{\theta(\psi)}_v(G_v)\).
\item Extend the assertions identified in the source as Theorem 1.5(2) to Archimedean places. In particular, for every such Archimedean \(v\), every permitted \(\psi\), and every pair of admissible auxiliary groups \(H_1,H_2\), establish
\[
\Pi^{\theta(\psi)}_{v,H_1}(G_v)
=
\Pi^{\theta(\psi)}_{v,H_2}(G_v).
\]
If \(G_v\) is quasi-split, establish, for every admissible \(H\),
\[
\Pi^{\theta(\psi)}_{v,H}(G_v)
=
\Pi^{\mathrm{Arthur}}_v(G_v,\psi),
\]
where the right-hand side denotes Arthur's packet for the corresponding parameter, as intended by the cited source.
\end{enumerate}

These equalities express the requested agreement of packets; the supplied record does not determine whether agreement includes multiplicities, equivalence conventions, or internal parametrizations. Likewise, no explicit endoscopic identity can be specified from the record alone: the relevant endoscopic data, character coefficients, transfer factors, normalizations, and parameter correspondence are omitted. The full statement of Theorem 1.5(2) and its hypotheses is also absent. Thus this is a source-dependent formulation requiring those missing specifications before it becomes a fully explicit mathematical statement.
\subsection{Short English statement}
For the local packets constructed from Arthur parameters in the cited source, establish the endoscopic character identities. Extend the stated auxiliary-group independence and quasi-split comparison with Arthur's packets to Archimedean places, retaining the hypotheses of Theorem 1.5(2). The supplied record omits the construction and several hypotheses needed for a fully explicit formulation.
\subsection{Sources}
\begin{itemize}
\item Supplied local UnsolvedMath registry record, permanent ID \(30004799\), source problem code \texttt{catalogue entry}, titled ``Endoscopic Identities and Archimedean Independence of Local Packets.'' The record belongs to the Oberwolfach Reports Open Problems collection, gives proposed year 2021, and names no individual proposer.
\item Supplied report citation: \emph{Automorphic Forms, Geometry and Arithmetic} (2022), Oberwolfach Reports, pp.~2089--2156, DOI \texttt{10.4171/owr/2021/39}. \url{https://doi.org/10.4171/owr/2021/39}. This is also the supplied literature source labeled ``original Oberwolfach report.'' The original statement refers to Theorem 1.5(2), whose text is not included in the record.
\item Supplied literature source: ``endoscopic character identity for even special orthogonal groups.'' \url{https://arxiv.org/abs/2506.17907}.
\item Supplied literature source: ``Chen--Zou, theta correspondence and Arthur packets.'' \url{https://arxiv.org/abs/2104.12354}.
\item Supplied dated literature curation, checked \texttt{2026-08-22}: the registry labels the problem ``partially solved'' and reports endoscopic character identities in important non-Archimedean cases, including recent results for bounded packets of non-quasi-split even orthogonal groups, while describing the extracted Archimedean and full local questions as unresolved. This assessment is supplied dated curation, not independently verified; its time-sensitive status and cited claims require independent verification.
\item The registry labels the statement \texttt{corrected\_verified} and states that the standalone statement and its hypotheses were checked against the official Oberwolfach report. This is a supplied verification claim, not an independent verification performed for this draft.
\end{itemize}
\end{document}
